\documentclass[a4paper,12pt]{extarticle}
\usepackage{geometry}
\usepackage{soul}
\sethlcolor{yellow}
\usepackage[T1]{fontenc}
\usepackage[utf8]{inputenc}
\usepackage[english,russian]{babel}
\usepackage{amsmath}
\usepackage{amsthm}
\usepackage{amssymb}
\usepackage{fancyhdr}
\usepackage{setspace}
\usepackage{graphicx}
\usepackage{colortbl}
\usepackage{tikz}
\usepackage{pgf}
\usepackage{subcaption}
\usepackage{listings}
\usepackage{indentfirst}
\usepackage[
backend=biber,
style=numeric,
maxbibnames=99
]{biblatex}
\addbibresource{refs.bib}
\usepackage[colorlinks,citecolor=blue,linkcolor=blue,bookmarks=false,hypertexnames=true, urlcolor=blue]{hyperref} 
\usepackage{indentfirst}
\usepackage{mathtools}
\usepackage{booktabs}
\usepackage[flushleft]{threeparttable}
\usepackage{tablefootnote}
\usepackage{chngcntr} % нумерация графиков и таблиц по секциям
\counterwithin{table}{section}
\counterwithin{figure}{section}
\usepackage[T2A]{fontenc}
\usepackage[utf8]{inputenc}
\usepackage[russian]{babel}
\usepackage{csquotes}
\usepackage{microtype} % для улучшения верстки

\graphicspath{{graphics/}}%путь к рисункам

\makeatletter
% \renewcommand{\@biblabel}[1]{#1.} % Заменяем библиографию с квадратных скобок на точку:
\makeatother

\geometry{left=2.5cm}% левое поле
\geometry{right=1.0cm}% правое поле
\geometry{top=2.0cm}% верхнее поле
\geometry{bottom=2.0cm}% нижнее поле
\setlength{\parindent}{1.25cm}
\renewcommand{\baselinestretch}{1.5} % междустрочный интервал


\newcommand{\bibref}[3]{\hyperlink{#1}{#2 (#3)}} % biblabel, authors, year
\addto\captionsrussian{\def\refname{Список литературы (или источников)}} 

\renewcommand{\theenumi}{\arabic{enumi}}% Меняем везде перечисления на цифра.цифра
\renewcommand{\labelenumi}{\arabic{enumi}}% Меняем везде перечисления на цифра.цифра
\renewcommand{\theenumii}{.\arabic{enumii}}% Меняем везде перечисления на цифра.цифра
\renewcommand{\labelenumii}{\arabic{enumi}.\arabic{enumii}.}% Меняем везде перечисления на цифра.цифра
\renewcommand{\theenumiii}{.\arabic{enumiii}}% Меняем везде перечисления на цифра.цифра
\renewcommand{\labelenumiii}{\arabic{enumi}.\arabic{enumii}.\arabic{enumiii}.}% Меняем везде перечисления на цифра.цифра

\begin{document}
\input{title_vkr}% это титульный лист - выберите подходящий вам из имеющихся в проекте вариантов (kr - курсовая работа у 3 курса, vkr - выпускная квалификационная работа у 4 курса)
\newpage
\setcounter{page}{2}

{
	\hypersetup{linkcolor=black}
	\tableofcontents
}

\newpage

\newpage
\section*{Аннотация}   % this is how to use russian
Ваша аннотация на русском языке.

\addcontentsline{toc}{section}{Аннотация}

\section*{Ключевые слова}
Рунические надписи, оптическое распознавание текста, OCR, TrOCR, мультимодальные модели, синтетические данные, малоресурсное обучение, доменный сдвиг, цифровая эпиграфика, древнескандинавский язык
\pagebreak

\section*{Введение}
\addcontentsline{toc}{section}{Введение}
 
Древнегерманские рунические надписи охватывают период приблизительно со II по XV в.\ н.\,э., локализованы преимущественно на территории Скандинавии, Британских островов и континентальной Европы и зафиксированы на камне, металле, кости и дереве. Несмотря на обширный научный опыт их изучения, значительная часть рунологической работы по-прежнему опирается на ручную транскрипцию специалистом, владеющим палеографической и археологической компетенциями. Это создаёт узкие места при оцифровке музейных коллекций, ограничивает воспроизводимость отдельных чтений и затрудняет доступ исследовательского сообщества к первичному материалу. В настоящей работе рунические надписи рассматриваются как объект автоматической обработки методами компьютерного зрения и обработки естественного языка.
 
\textbf{Актуальность исследования} обусловлена тремя обстоятельствами. Во-первых, в 2019--2025 гг.\ методы машинного обучения закрепились в обработке древних текстов: модели Pythia, Ithaca и Aeneas установили стандарт для восстановления и атрибуции древнегреческих и латинских эпиграфических текстов~\cite{assael2019,assael2022,assael2025}, а обзор в журнале \textit{Computational Linguistics} систематизировал направление как самостоятельную ветвь вычислительной филологии~\cite{sommerschield2023}. Во-вторых, сохраняется разрыв между объёмом рунического корпуса и доступными цифровыми возможностями: значительная часть скандинавских архивов оцифрована неполно, а экспертная транскрипция остаётся трудоёмкой~\cite{magin2024}. В-третьих, трансформерные архитектуры оптического распознавания типа TrOCR демонстрируют устойчивое преимущество над свёрточно-рекуррентными конвейерами на задачах исторического OCR~\cite{li2023}, а методы параметрически эффективной адаптации (PEFT, в частности LoRA) позволяют переносить крупные предобученные модели в малоресурсные домены при ограниченных вычислительных бюджетах~\cite{hu2022}.
 
\textbf{Степень разработанности проблемы} характеризуется выраженной асимметрией между эпиграфическими традициями. Для греко-латинской традиции сложился методологический стандарт линейки Pythia~-- Ithaca~-- Aeneas~\cite{assael2019,assael2022,assael2025}. Клинописная эпиграфика обработана как на уровне распознавания и транслитерации знаков~\cite{gordin2020}, так и на уровне восстановления повреждённых фрагментов~\cite{fetaya2020}. Конвейер DeepScribe демонстрирует применимость модульных систем компьютерного зрения к корпусу эламской клинописи~\cite{williams2025}. Для китайской традиции цзягувэнь разработаны двухступенчатые системы детекции и классификации знаков~\cite{fujikawa2023} и генеративные подходы к дешифровке средствами диффузионных моделей~\cite{guan2024}. Близкая по форме древнетюркская руниформа получила конвейеры синтетической генерации данных с трёхмерным моделированием носителя~\cite{derin2024}, нейросетевые модели для орхонских надписей~\cite{mukashev2024} и базовые архитектуры распознавания~\cite{taheri2025}. Двухступенчатая схема «детекция + классификация» проработана также применительно к греческим надписям на византийских печатях~\cite{rageau2025}. Вместе с тем на момент первой половины 2026 г.\ в индексируемых рецензируемых изданиях отсутствуют публикации, в которых трансформерные OCR-архитектуры применялись бы к корпусу скандинавских рунических надписей. Тем самым формируется локализованная исследовательская лакуна, заполнение которой составляет содержание настоящей работы.
 
\textbf{Объектом исследования} выступают памятники древнегерманской рунической эпиграфики на различных типах материальных носителей. \textbf{Предметом исследования} являются методы и алгоритмы автоматического распознавания рунических знаков на фотографических изображениях в условиях малоресурсного домена.
 
\textbf{Цель работы}~-- разработать и экспериментально оценить систему
автоматического распознавания и лингвистического анализа древнегерманских
рунических надписей по фотографическим изображениям с применением современных
методов глубокого обучения и синтетической аугментации данных. Центральный исследовательский вопрос: в какой мере комбинация контекстных трансформерных архитектур и контролируемой синтетической генерации данных позволяет преодолеть малоресурсность рунического материала и приблизить долю посимвольных ошибок
(Character Error Rate, CER)  к условному порогу практической применимости -- 20\,\%? \hl{Под лингвистическим анализом здесь понимается
демонстрационный модуль интерпретации распознанной транслитерации, опирающийся
на существующие лексические и нормализационные ресурсы, а не система
автоматического перевода, построение которой выходит за рамки работы.}
 
Для достижения цели решаются следующие \textbf{задачи}:
\begin{enumerate}
    \item провести систематический обзор современных подходов компьютерного зрения и машинного обучения к распознаванию (OCR/HTR) исторических и эпиграфических памятников;
    \item сформировать эмпирическую базу: собрать и систематизировать корпус фотографий реальных рунических надписей из открытых архивов и разработать алгоритм синтеза искусственных изображений для расширения обучающей выборки;
    \item спроектировать и реализовать конвейер предобработки изображений, адаптированный к физической и спектральной деградации эпиграфических носителей;
        \item Реализовать и сравнить два подхода к сквозному (end-to-end) распознаванию транслитерации: дообучение специализированной OCR-модели TrOCR и адаптацию мультимодальной vision-language-модели семейства Qwen-VL с использованием метода параметрически эффективного дообучения LoRA на синтетическом корпусе данных;
    \item выполнить абляционное исследование (ablation study) для количественной оценки вклада синтетической аугментации, предобработки и выбора энкодера в итоговое качество;
    \item \hl{разработать демонстрационный модуль автоматической лингвистической интерпретации распознанной транслитерации;}
    \item провести качественный и количественный анализ типологии ошибок и определить границы практической применимости системы.
\end{enumerate}
 
\textbf{Методология} сочетает теоретические, эмпирические и экспериментальные процедуры: аналитический обзор литературы и типологизацию архитектур; отбор, очистку и аннотацию изображений из открытых архивов и формирование контролируемого синтетического корпуса; обучение и дообучение моделей с использованием transfer learning, сравнение архитектурных конфигураций и абляционные эксперименты по ключевым компонентам. Ввиду отсутствия унифицированного изобразительно-текстового корпуса рунических
надписей обучение моделей ведётся
на синтетическом корпусе, а оценка~-- на вручную выверенном реальном «золотом»
множестве ограниченного объёма, что изолирует и количественно характеризует
синтетически-реальный доменный сдвиг.
 
\textbf{Эмпирическую базу} составляют четыре источника: Samnordisk runtextdatabas (Scandinavian Runic-text Database, SRD) Уппсальского университета, предоставляющая транскрипции и нормализованные тексты; поисковая служба Söktjänsten Runor Шведского управления по охране культурного наследия (Riksantikvarieämbetet), агрегирующая метаданные около семи тысяч надписей; открытые архивы скандинавских музеев (Historiska museet, Danske Runeindskrifter, Runearkivet Музея истории культуры Университета Осло); синтетический корпус собственной генерации с моделью физической деградации носителя.
 
\textbf{Научная новизна} работы состоит в следующем: представлена первая систематическая адаптация трансформерных OCR-архитектур к скандинавскому руническому материалу; предложена методология контролируемой синтетической генерации рунических изображений с явным моделированием деградации носителя; выполнено сравнительное сопоставление OCR-архитектуры TrOCR
и мультимодальной vision-language-модели (Qwen-VL), адаптированных параметрически
эффективным методом (LoRA) на синтетическом корпусе, в рамках единого протокола оценки; разработана унифицированная схема аннотации, совместимая с задачами компьютерного зрения и допускающая повторное использование.
 
\textbf{Теоретическая значимость} определяется первым систематическим тестированием трансформерных OCR-архитектур на корпусе скандинавских рунических надписей, что расширяет эмпирическую базу обзорных исследований по машинному обучению для древних языков, вкладом в методологию переноса техник синтетического обучения между эпиграфическими доменами и обоснованием эмпирических порогов CER для low-resource эпиграфики. \textbf{Практическая значимость} состоит в формировании воспроизводимого инструментария (размеченного датасета, исходного кода конвейера и весов обученных моделей) и в потенциальном ускорении оцифровки архивных коллекций как основы для последующих задач атрибуции, датировки и реконструкции лакун.
 
\textbf{Структура работы} следует логике от обзора и теоретических оснований к экспериментальной верификации. Первая глава содержит критический обзор литературы по машинному обучению для древних языков и эпиграфики. Вторая глава излагает рунологические, лингвистические и архитектурно-вычислительные основания работы. Третья глава описывает эмпирическую базу, процедуры аннотации и предобработки изображений с учётом деградации носителя. Четвёртая глава представляет методологию и архитектуру разработанного
конвейера~-- две end-to-end-конфигурации (дообучение TrOCR и адаптацию
мультимодальной модели Qwen-VL средствами LoRA), а также демонстрационный
модуль лингвистической интерпретации транслитерации. Пятая глава содержит экспериментальную часть: протокол оценки, основные результаты, ablation study и типологический анализ ошибок. Заключение подводит итоги и формулирует перспективы исследования. Работу завершает список использованной литературы.




\section{Обзор литературы}
\label{sec:litreview}

Настоящая глава представляет критический обзор работ 2019--2026 гг.\ по автоматическому распознаванию (OCR/HTR) исторических текстов на твёрдых носителях средствами глубокого обучения. Цель обзора~-- очертить состояние разработанности темы и локализовать незаполненную нишу. Центральный результат, к которому подводит анализ, состоит в следующем: по состоянию на 2026 г.\ в рецензируемой литературе отсутствуют публикации, применяющие современные нейросетевые архитектуры (CNN, ViT, TrOCR, мультимодальные модели) к надписям Старшим, Младшим или Англосаксонским футарком. Немногочисленные смежные работы относятся либо к древнетюркскому руническому письму, либо к традиции 3D-сканирования рунических камней. Этот пробел определяет структуру обзора: сначала в работе рассматриваются семейства архитектур (трансформерный HTR, TrOCR, мультимодальные модели, детекторы объектов), затем~-- рунический материал и цифровая рунология, а также методологические аналогии из смежных малоресурсных эпиграфических доменов.

\hl{Поиск проводился в репозиториях arXiv, ACL Anthology, IEEE Xplore, Springer, ACM Digital Library и базах Nature и PLOS, а также в специализированных журналах \textit{Futhark: International Journal of Runic Studies}, ACM \textit{Journal on Computing and Cultural Heritage} (JOCCH), \textit{International Journal on Document Analysis and Recognition} (IJDAR) и \textit{Computational Linguistics}. Дополнительно просматривались порталы проектов RuneS, Rundata и Riksantikvarieämbetet.}

\subsection{Трансформерные архитектуры в распознавании рукописного текста}
\label{sec:lit-htr}

В период 2020--2022 гг.\ трансформерные архитектуры начали вытеснять традиционные конвейеры CNN\,+\,BiLSTM\,+\,CTC в задачах распознавания рукописного текста. Kang и др.~\cite{kang2022} предложили одну из первых полностью трансформерных архитектур HTR, объединившую визуальный экстрактор ResNet-50 с мультиголовым самовниманием на визуальной и текстовой сторонах, и достигли на корпусе IAM CER около 4,67\,\% при обучении на синтетических строках. Diaz и др.~\cite{diaz2021} провели систематическое сравнение энкодер-декодерных пар и показали, что Transformer-энкодер в роли последовательностного экстрактора превосходит рекуррентные альтернативы, тогда как авторегрессивный Transformer-декодер не даёт прироста на коротких строковых задачах. Переход к страничному распознаванию совершила архитектура DAN (Document Attention Network) Coquenet и др.~\cite{coquenet2023}, объединившая FCN-энкодер с трансформерным декодером, генерирующим символы вместе с тегами логической структуры.

Наиболее показательна архитектура HTR-VT~\cite{li2025htrvt}~-- чисто энкодерный ViT с CNN-эмбеддингом, обучаемый без внешних данных: он достигает CER 4,7\,\% на IAM и 3,9\,\% на READ~2016. Авторы особенно отмечают важность абляции: удаление CNN-фронтенда поднимает CER на IAM с 3,3\,\% до 26,6\,\%, что опровергает представление о «чистом» трансформере как универсальном решении для распознавания письма и обосновывает гибридную CNN\,+\,Transformer-схему как методологическую норму. Однако эффективность применения этого метода для нелатинской письменности ещё не освещалась; ни одной публикации, применяющей ViT-семейство к руническим скриптам в полноценном HTR-режиме, обнаружить не удалось.

\subsection{TrOCR и адаптация к историческим документам}
\label{sec:lit-trocr}

Модель TrOCR~\cite{li2023} стала стандартом современных исследований распознавания исторического письма. Её архитектура объединяет энкодер типа ViT/BEiT/DeiT с декодером типа RoBERTa/MiniLM и обучается на двухступенчатом конвейере из сотен миллионов синтетических печатных и миллионов рукописных строк, достигая на IAM CER 2,89\,\% (вариант LARGE), 3,42\,\% (BASE) и 4,22\,\% (SMALL). Ценность TrOCR для эпиграфики определяется не столько абсолютными результатами, сколько возможностью дообучения относительно небольшим числом размеченных строк. Этот вопрос систематически изучен Ströbel и др.~\cite{strobel2022}: дообученный TrOCR-LARGE превосходит специализированную систему Transkribus HTR+ на исторических письмах XVI~в.\ уже после пятнадцати эпох, не используя при этом информации о базовых линиях, необходимой Transkribus. Последующая работа~\cite{strobel2023} уточняет, что устойчивое превосходство над свёрточно-рекуррентными моделями демонстрирует только вариант LARGE.

Систематическое продолжение представлено у Parres и Paredes~\cite{parres2023}, где конфигурация Vision\,Encoder--Decoder даёт WER 6,9\,\% на ICFHR~2014 Bentham и 14,5\,\% на ICFHR~2016 Ratsprotokolle; авторы показывают, что при кросс-языковой адаптации необходимо дообучать и декодер, а его замораживание ведёт к потере 3--5~п.п.\ CER. Декодер-only модель DTrOCR~\cite{fujitake2024} доводит IAM до CER 2,38\,\%, формируя новый ориентир. Адаптация TrOCR к средневековым рукописям закреплена в проекте CATMuS Medieval~\cite{clerice2024}. Самообучаемое предобучение трансформерных распознавателей развивают работы Souibgui и др.~\cite{souibgui2023textdiae}, применившие masked-image-modelling с реконструкцией деградации, что особенно ценно для деградированных каменных носителей, и Vogler и др.~\cite{vogler2022}, показавшие, что тридцати размеченных строк достаточно для адаптации MAE-предобученного энкодера к новой рукописной традиции~-- это наиболее прямой методологический ориентир для рунической задачи. Применений TrOCR к руническим надписям не существует.

\subsection{Мультимодальные модели в OCR исторических документов}
\label{sec:lit-vlm}

Распространение визуально-языковых моделей в 2023--2025 гг.\ породило волну попыток оценить их пригодность для исторического OCR. Архитектурный ландшафт задают LLaVA~\cite{liu2023llava}, семейство Qwen-VL~\cite{wang2024qwen2vl}, InternVL~\cite{chen2024internvl}, компактная MiniCPM-V~\cite{yao2025minicpm} и Florence-2~\cite{xiao2024florence}, чьи способность к OCR-with-region и открытая лицензия делают её привлекательным кандидатом для дообучения на эпиграфических корпусах; закрытые GPT-4V/4o и Gemini~\cite{geminiteam2023} остаются бенчмарками для zero/few-shot-экспериментов. Систематическая оценка OCR-способностей этих моделей закреплена в OCRBench~\cite{liu2024ocrbench} и его расширении OCRBench~v2~\cite{fu2025ocrbench2}: 36 из 38 оценённых моделей получили менее 50 из 100 баллов, причём систематические провалы приходятся на редкие скрипты, мелкие символы и нестандартную композицию~-- все три характеристики применимы к рунической эпиграфике.

Для древних языков ориентиром остаются работы группы DeepMind/Oxford: Pythia~\cite{assael2019}, Ithaca~\cite{assael2022} и мультимодальная Aeneas~\cite{assael2025}, выполняющие реставрацию лакун, географическую атрибуцию и датировку греческих и латинских надписей. Cowen-Breen и др.~\cite{cowenbreen2024} предлагают instruction-tuned causal LM, конкурирующий с Ithaca, что отражает миграцию методологии в сторону декодер-only моделей. Систематических оценок мультимодальных моделей на руническом материале не обнаружено.

\subsection{Детекторы объектов для эпиграфики}
\label{sec:lit-detect}

В типичном эпиграфическом конвейере детекция строк или знаков предшествует распознаванию, поэтому ландшафт детекторов остаётся существенным. Семейство YOLO~\cite{jocher2020,wang2024yolov9} прошло путь от ранних версий без формальной публикации до архитектур с программируемым градиентом и NMS-free-обучением. При этом сравнительные оценки на нестандартных доменах показывают, что новизна версии не гарантирует превосходства на узких исторических данных. Специализированные текстовые детекторы CRAFT~\cite{baek2019craft} и DBNet~\cite{liao2020dbnet} остаются обязательным компонентом эпиграфических конвейеров, поскольку строки на каменных носителях рутинно искривлены и повёрнуты.

Особенно важна для рунических надписей ветвь ориентированных ограничивающих рамок (Oriented Bounding Box, OBB): рунические строки на стелах нередко идут под произвольными углами, по дугам или спиралям. Ключевые ориентиры~-- Oriented R-CNN~\cite{xie2021orientedrcnn} и ReDet~\cite{han2021redet}, а также OBB-варианты детекторов YOLO. Обзор сегментации строк и детекции базовых линий в исторических документах~\cite{brodic2025} и работа YOLO-HTR~\cite{lomov2025}, интегрирующая YOLO-детектор с Vertical Attention Network, систематизируют методическую базу. Прямых применений OBB-семейства к руническим надписям не опубликовано, что образует ещё одну отчётливую лакуну.

Тем не менее, специфика рунического материала вносит существенные коррективы в классическую каскадную схему компьютерного зрения. Полное отсутствие публично доступных размеченных корпусов с посимвольной или строчной локализацией (bounding boxes), а также критически малый объем дошедших до нас физических артефактов делают обучение робастных OBB-детекторов «с нуля» практически нереализуемым на данном этапе. Более того, высокая степень эрозии и фрагментарности носителей приводит к тому, что границы отдельных знаков или строк оказываются семантически размытыми даже для экспертов-рунологов. По этой причине в рамках данного исследования обосновывается отказ от традиционного двухэтапного конвейера «детекция → классификация» в пользу сквозного (end-to-end) распознавания целых текстовых полей методами мультимодальных трансформеров. 

\subsection{Рунические надписи: цифровая рунология и состояние корпуса}
\label{sec:lit-runic}


\subsubsection{Корпуса и инфраструктура}

Стержневым инфраструктурным проектом для германских рун является RuneS (\textit{Runic Writing in the Germanic Languages})~\cite{runesdb} при Союзе академий наук Германии, охватывающий руническое письмо примерно AD~100--1500 и агрегирующий данные Кильского рунического проекта, датского регистра и Скандинавской руно-текстовой базы под открытыми лицензиями. RuneS, однако, представляет собой структурированную базу и графемно-типологическую платформу, а не среду машинного обучения; связанных с ним нейросетевых публикаций не найдено. Параллельная скандинавская инфраструктура опирается на Rundata / Samnordisk runtextdatabas~\cite{bianchi_rundata}; несмотря на статус де-факто стандартного транскрипционного корпуса, в рецензируемой литературе не выявлено ни одной работы, где Rundata использовалась бы как обучающий или валидационный материал для машинного обучения. Программа \textit{Everlasting Runes} и обзорная статья Källström~\cite{kallstrom2022} о корпусных изданиях шведских рунических надписей также носят характер цифрового документирования без ML-методологии.

\subsubsection{Традиция 3D-сканирования}

Наиболее зрелым вычислительным направлением для германских рун является не глубокое обучение, а анализ структуры поверхности средствами 3D-лазерного сканирования. Наиболее цитируемая работа Imer, Kitzler Åhfeldt и Zedig~\cite{imer2023} показывает, что 3D-сканирование в сочетании с PCA позволяет идентифицировать «руку» мастера и связать камни Йеллинга с конкретным резчиком. Авторы признают ограничения подхода~-- малый размер выборки, субъективность интерпретации, шум выветривания, что само по себе мотивирует переход к нейросетевым методам метрического обучения, пока никем не предпринятый. Группа открыла 3D-датасеты, на которые легли в основу их исследования, однако предложенные наработки не применимы в контексте данной задачи.

\subsubsection{Состояние цифровой рунологии и единичные нейросетевые работы}

Анализ работ последних опубликованных томов журнала \textit{Futhark} (2020--2024) ясно показывает: ни одна статья не применяет машинное обучение или компьютерное зрение к руническим надписям, что может говорить о малом освещении использования методов ML для «рунических» скандинавских алфавитов. Реальный объём нейросетевых публикаций по «руническим» скриптам сосредоточен в области древнетюркской руники, названной «рунической» в силу визуального сходства с германскими футарками, но лингвистически восходящей к арамейскому письму. Эти работы остаются ближайшими методологическими ориентирами. Mukashev и Tukeyev~\cite{mukashev2024} описывают нейросетевой конвейер распознавания орхонских рун, хотя и полное описание архитектуры, а также метрик не удалось найти в открытом доступе. Taheri и др.~\cite{taheri2025} обучают компактную CNN на 38 классах древнетюркских символов и сообщают точность 96,3\,\%, признавая ограничения~-- крохотный исходный датасет и потерю деталей штриха. Методически наиболее изощрённой является работа Derin и Uçar~\cite{derin2024}, где Vision Transformer обучается на полностью синтетических данных, полученных через параметрическую рандомизацию и физически-основанный рендеринг в Blender, и применяется в zero-shot-режиме к фотографиям камня Кюль-Тегина с высокой точностью без единого реального обучающего примера. Этот результат~-- наиболее прямой ориентир для настоящей работы: он демонстрирует, что синтетический рендеринг каменной поверхности с параметризованной деформацией графем способен заменить размеченный корпус, что прямо отвечает на главное ограничение рунической задачи.

\subsubsection{Германские руны: отсутствие рецензируемых работ}

Применительно к собственно германским футаркам рецензируемые нейросетевые публикации отсутствуют. На IEEE DataPort размещён датасет рукописного Старшего футарка (25 классов) без сопровождающей рецензируемой публикации. Просмотр скандинавских репозиториев диссертаций (DiVA, Munin) не выявил ни одной работы, применяющей глубокое обучение или компьютерное зрение к германским руническим надписям; имеющиеся рунологические диссертации используют стереомикроскоп и классическую палеографию. Это прямая лакуна, образующая исследовательскую нишу настоящей работы.

\subsection{Методологические аналогии из смежных малоресурсных доменов}
\label{sec:lit-analog}

Поскольку прямых аналогов в рунологии немного, методологический инструментарий заимствуется из родственных малоресурсных областей. Наиболее проработанный пример~-- клинопись. Dencker и др.~\cite{dencker2020} вводят слабосупервизорное итеративное обучение детектора знаков на основе выравнивания латинских транслитераций с изображениями, получая 8109 размеченных знаков 186 классов без дорогой плотной разметки; эта идея напрямую переносима на рунический материал, где транслитерации Rundata могут выступить тем же выравнивающим сигналом. Гейдельбергская группа~\cite{bogacz2022} развивает 3D-подход, приводящий задачу к квазидвумерной форме через анализ кривизны поверхности.

Отдельную методическую линию образует синтетическая генерация изображений, центральная для малоресурсной задачи. Базовые диффузионные модели~\cite{rombach2022ldm} и условный контроль через ControlNet~\cite{zhang2023controlnet} обеспечивают управляемый синтез даже при ограниченных обучающих данных; линия GAN-синтеза рукописного письма представлена ScrabbleGAN~\cite{fogel2020scrabblegan}. Также, для условий малой выборки релевантны канонические методы few-shot- и self-supervised-обучения.

\subsection{Идентифицированные лакуны и позиционирование работы}
\label{sec:lit-gaps}

Проведённый обзор позволяет зафиксировать ряд лакун, образующих обоснование настоящей работы. Во-первых, отсутствует нейросетевая модель распознавания, обученная или адаптированная под Старший, Младший или Англосаксонский футарк, тогда как смежные традиции~-- клинописная, оракулькостная, греко-латинская, египтологическая, майяская~-- уже располагают собственными ML-направлениями. Во-вторых, существующие скандинавские 3D-корпуса анализируются исключительно классической многомерной статистикой, без применения нейросетевых архитектур, оперирующих трёхмерными данными. В-третьих, автоматическая сегментация рунических знаков на фотографиях не реализована: ни детекторы YOLO/DETR, ни OBB-семейство к корпусу германских рунических камней в опубликованной форме не применялись. В-четвёртых, синтетическая генерация рунических изображений с физически-основанным рендерингом и моделированием выветривания не разработана, хотя прямые шаблоны для близких доменов уже существуют. Наконец, систематические оценки мультимодальных моделей на руническом материале отсутствуют.

Таким образом, современная литература по распознаванию исторических эпиграфических надписей вышла на технологический уровень, существенно превышающий уровень любых опубликованных работ по руническому материалу. Настоящая работа позиционируется как систематический перенос методологии из перечисленных смежных областей на корпус скандинавских рунических надписей, опирающийся на инфраструктуру RuneS и Rundata и на методы синтетической генерации.

\section{Лингвистические и материальные основания задачи}
\label{sec:foundations}


\subsection{Рунические письменные системы: графемный инвентарь и история}
\label{sec:writing-systems}
 
В отличие от типичной задачи OCR, где разнообразие исходных алфавитов покрывается единой моделью на уровне Unicode-кодов, рунический материал распадается на три исторически различные системы~-- Elder Futhark, Younger Futhark и Anglo-Saxon Futhorc, каждая со своим графемным инвентарём, своим соотношением знака и фонемы и своим репертуаром носителей. Эти различия прямо переводятся в инженерные параметры задачи: размер пространства классов, степень омонимии знака, ожидаемый уровень доменного сдвига между сегментами выборки и допустимую структуру целевой метрики. Поэтому каждое свойство системы соотносится ниже с конкретной трудностью, которую она создаёт для классификатора графем и последующего лингвистического слоя.
 
\subsubsection{Старший футарк}
 
Старший футарк (Elder Futhark)~-- древнейшая засвидетельствованная германская рунописная система, фиксируемая на территории от южной Скандинавии до Карпатского бассейна приблизительно со II по VIII в.\ н.\,э., с пиком документации между серединой III и серединой VII в.~\cite{looijenga2003,duwel2008,barnes2012}. Канонический инвентарь состоит из двадцати четырёх знаков, организованных в три группы по восемь рун (\textit{ættir}); порядок групп и последовательность знаков сохраняются во всех известных полных рядах, начиная с надписи на камне из Килвера (Готланд, около 400~г.\ н.\,э.), где система впервые засвидетельствована \textit{in extenso} в каноническом порядке \texttt{f u \Thorn{} a r k g w | h n i j \"i p z s | t b e m l \Eng{} d o}~\cite{krause1966,antonsen1975}.
 
Лингвистическим основанием системы считается принцип «один знак~-- одна фонема»: набор графем приблизительно соответствует фонемному инвентарю прагерманского и раннеобщескандинавского~\cite{antonsen1975,looijenga2003}. Однако, этот принцип может и нарушаться. Графема \textipa{ï} не имеет устойчивого фонетического значения и в разных интерпретациях возводится к долгому переднему гласному либо к специфическому рефлексу прагерманского, что отражает реальную неопределённость прочтения~\cite{krause1966,barnes2012}. Систематически фиксируются связные руны (\textit{bind-runes}) и зеркальные начертания, обусловленные эстетикой носителя и техническими ограничениями резца~\cite{looijenga2003}. Тем не менее степень омонимии графем здесь относительно низка, и в нормативной транслитерации большинство знаков отображается на единственную фонему.
 
Корпус Старшего футарка географически и материально разнороден, что определяет визуальную гетерогенность доступной выборки: помимо камня из Килвера, опорными памятниками являются железные наконечники копий из Эвре-Стабу и находки из болота Иллеруп-Одаль, а существенный сегмент образуют золотые брактеаты V--VI~вв.~\cite{looijenga2003,duwel2008}. Таким образом, система документирована на камне, железе, драгоценных металлах, кости и органике, причём пропорции категорий в сохранившемся корпусе определяются факторами археологической сохранности, а не первоначальным распределением употреблений~\cite{looijenga2003,barnes2012}.
 
Для распознавания отсюда следует двойственный вывод. Относительная фонематическая прозрачность делает классификацию графемы внутри этой системы концептуально более чистой, чем в Младшем футарке: верно опознанный знак в большинстве случаев соответствует единственной фонеме, и ошибки CER близки к ошибкам на фонемном уровне. Однако материальная и хронологическая разнородность корпуса формирует значительный внутрисистемный доменный сдвиг: модель, обученная преимущественно на камне, не переносится на металлические подложки без доменной адаптации. Малочисленность сохранившихся надписей превращает Старший футарк в типичный малоресурсный сегмент, где обучение без few-shot learning и синтетической аугментации невозможно.
 
\subsubsection{Младший футарк}
 
Переход к Младшему футарку (Younger Futhark), завершившийся к концу VIII в.~\cite{schulte2011,stoklund2010}, образует одну из наиболее обсуждаемых проблем рунологии и узловой пункт для постановки задачи фонемного распознавания. Суть проблемы~-- парадокс редукции: число графем сокращается с двадцати четырёх до шестнадцати, тогда как фонологическая система древнескандинавского языка в тот же период усложняется (формируются новые гласные качества как результат i- и u-перегласовки, развиваются длительностные противопоставления, дифференцируются носовые)~\cite{schulte2011,barnes2012}. Шульте описывает это движение как «невидимую руку перемен», где функциональная редукция инвентаря сопровождается, а не компенсируется, ростом фонологической сложности~\cite{schulte2011}.
 
Механизм адаптации состоит в нейтрализующем слиянии противопоставлений на письме при сохранении их в речи. Противопоставления звонкости и глухости в смычных рядах \textipa{/p/--/b/}, \textipa{/t/--/d/}, \textipa{/k/--/g/} исчезают с графемного уровня: одна руна обслуживает обе фонемы. Графема \textit{kaun} функционирует как письменное выражение по меньшей мере трёх единиц~-- \textipa{/k/}, \textipa{/g/} и носового \textipa{/ŋ/} перед велярным смычным, а \textit{týr} передаёт и \textipa{/t/}, и \textipa{/d/}~\cite{duwel2008,barnes2012,spurkland2005}. Знак и фонема находятся в отношении «один-ко-многим», и разрешение неоднозначности возможно лишь на уровне лексического и контекстного знания.
 
Отсюда следствие для метрик качества. Если классификатор правильно распознал визуальный знак \textit{kaun}, метрика CER на уровне графем фиксирует успех; однако при выходе на фонологическую транскрипцию та же графема может соответствовать \textipa{/k/}, \textipa{/g/} или \textipa{/ŋ/}, и фонемная ошибка возникнет независимо от качества визуальной классификации. Таким образом, графемный CER в Младшем футарке систематически не эквивалентен фонемному.
 
Внутренняя организация Младшего футарка неоднородна: уже в IX в.\ выделяются три региональные ветви. Длинноветвистые (long-branch, датские) руны представлены полными вертикалями и развитыми наклонными ветвями; коротковетвистые (short-twig, шведско-норвежские) редуцируют ветви до коротких штрихов; хэльсингские (staveless) руны удаляют у ряда знаков основной ствол~\cite{barnes2012,spurkland2005,duwel2008}. С точки зрения компьютерного зрения три ветви образуют три различных пространства аллографов одного фонемного инвентаря, что значимо для построения system-aware заголовка классификатора. Количественная плотность корпуса значительно превосходит плотность Старшего: число рунических камней с надписями Младшего футарка в одной только Швеции оценивается приблизительно в 1700--2500 объектов. Это означает, что в Rundata Младший футарк доминирует и формирует значительный class imbalance относительно остального материала. Наиболее протяжённой надписью является камень из Рёка (около 800~г.\ н.\,э.), включающий, по разным подсчётам, более семисот рун и одновременно использующий short-twig руны, отдельные знаки Старшего футарка и шифровальные системы~\cite{barnes2012}.
 
\subsubsection{Англосаксонский футорк}
 
Англосаксонский футорк (Anglo-Saxon Futhorc) развивается на острове и на фризской периферии в противоположном направлении: вместо редукции инвентарь расширяется под действием фонологических сдвигов древнеанглийского. Стандартный ряд включает двадцать восемь знаков; рукописные ряды (Cotton Domitian~A.IX и др.) фиксируют расширение до тридцати трёх рун~\cite{page1999,parsons1999,bammesberger1991}. Позднейший пласт (\textit{cweor\Thorn{}}, \textit{calc}, \textit{gar}, \textit{stan}), документированный почти исключительно в рукописях, Пейдж характеризует частично как «псевдо-руны», изобретённые для соответствия латинскому алфавиту~\cite{page1999,parsons1999}. Для классификатора это означает, что пространство классов не фиксировано априорно и зависит от состава обучающей выборки: модель, обученная только на каменной эпиграфике, может не встретить \textit{cweor\Thorn{}} и \textit{stan}.
 
 
\subsubsection{Следствия для архитектурных решений}
 
Сопоставление трёх систем даёт четыре вывода для проектирования модели. Во-первых, размер классификационного пространства варьируется от шестнадцати классов (Младший футарк) до приблизительно тридцати трёх (расширенный футорк) при крайне неравномерном распределении примеров, что требует взвешенной функции потерь либо стратифицированной выборки. Во-вторых, омонимия в Младшем футарке создаёт неустранимую неопределённость ground truth на уровне фонологической транскрипции, и оценка модели должна учитывать это через раздельные метрики CER (графемный) и PER (фонемный). В-третьих, поскольку правила соответствия графемы и фонемы зависят от системы, классификатор должен быть условен по идентификатору системы~-- требуется system-aware classifier head, переключаемый априорным предсказанием системы.
 
Таблица~\ref{tab:runic-systems} систематизирует ключевые параметры трёх систем в форме, пригодной для прямого соотнесения с инженерными решениями.
 
\begin{table}[htbp]
\centering
\caption{Сравнительная характеристика рунических письменных систем}
\label{tab:runic-systems}
\small
\begin{tabular}{p{3.0cm} p{3.1cm} p{3.1cm} p{3.4cm}}
\hline
\textbf{Параметр} & \textbf{Elder Futhark} & \textbf{Younger Futhark} & \textbf{Anglo-Saxon Futhorc} \\
\hline
Период & ок.\ 150--700 н.\,э. & ок.\ 800--1100 н.\,э. & ок.\ 650--1100 н.\,э. \\
Кол-во знаков & 24 & 16 & 28 (станд.)\,--\,33 (расш.) \\
Знаков на фонему (макс.) & 1 (с редкими искл.) & до 3 (\textit{kaun} = /k/, /g/, /\Eng/) & 1--2 (палатализация /k/, /g/) \\
Носители & камень, металл, кость, дерево (фрагм.) & камень (преим.), металл, кость, дерево & камень, кость, металл, пергамен (\textit{runica manuscripta}) \\
Лигатуры & регулярно & реже, засвидетельствованы & регулярно, особенно в рукописях \\
Следствие для CV/OCR & малый корпус, высокий доменный сдвиг по материалам & class imbalance в Rundata, графемно-фонемный разрыв & максимальный диамедиальный сдвиг (камень $\leftrightarrow$ пергамен) \\
\hline
\end{tabular}
\end{table}
 
Каждая система задаёт собственный набор инженерных требований, и попытка построить унифицированный классификатор без явного представления о системе будет страдать от усреднения разнородных распределений.
 
\subsection{Физическая деградация носителей как задача компьютерного зрения}
\label{sec:degradation}
 
Переход от лингвистического описания к материальному воплощению вводит второе принципиальное отличие рунической задачи от стандартного OCR. В документном OCR физическая деградация (пятна, разводы чернил, искажения сканирования) представляет собой технический шум, аддитивный по отношению к чистому сигналу и в значительной степени удалимый классическими методами. В эпиграфическом домене химические и биологические процессы изменяют именно те признаки~-- локальный контраст, текстурный градиент, цветовую вариативность,~-- на которых классические алгоритмы основывают извлечение сигнала. Деградация рунического носителя является, таким образом, не аддитивным, а структурным шумом, визуально неотличимым от сигнала средствами классического компьютерного зрения. Это и обосновывает отказ от Otsu, Sauvola и Canny как основных (а не вспомогательных) методов извлечения знака и переход к нейросетевому подходу. Объёмные методы документации (3D-сканирование, RTI) в работе не используются и упоминаются лишь как смежное направление.
 
\subsubsection{Типология носителей и механизмы деградации}
 
Камень составляет численно доминирующую категорию корпуса и наиболее уязвим для долгосрочных погодных и физико-химических воздействий. На большинстве снимков видны результаты воздействия лишайников и мхов на поверхность: под их воздействием разрушаются исходные текстуры, что приводит к размытию границ рун. Параллельно идёт колонизация мхами, усиливающими разрушение в циклах замерзания--оттаивания. 

Деградация металлических носителей происходит по-разному: в зонах антропогенного реза коррозия развивается иначе, чем на нетронутой поверхности, благодаря этому есть возможность отделить намеренный рез от случайного повреждения носителя.
 
Кость и рог, представленные мелкими артефактами, подвержены дегидратации и денатурации коллагена. Продольное растрескивание создаёт паттерн линейных разрывов, визуально близкий к резу и образующий ещё один источник ложноположительных детекций. Дерево как носитель в европейском климате почти полностью утрачено. Деревянный сегмент географически локализован и обладает специфическим визуальным режимом (продольная древесная текстура с тёмными резами), не моделируемым аугментациями каменного материала.
 
\subsubsection{Спектральная неоднозначность}
 
Механизмы деградации сходятся в одной проблеме компьютерного зрения: антропогенный рез и естественное повреждение на двумерной фотографии обладают почти идентичными характеристиками локального цвета, контраста и градиента яркости. В пространстве низкоуровневых дескрипторов класс «руническая линия» неотделим от класса «естественная неоднородность поверхности».

Условия съёмки влияют на эту неотделимость сильнее, чем выбор алгоритма обработки. Боковое касательное освещение подчёркивает рельеф и делает рез различимым, но одновременно усиливает все естественные неровности~\cite{earl2010,calia2023}. Диффузное освещение нивелирует рельеф и удаляет тени, улучшая цветовую однородность, но лишая изображение информации о глубине реза. Различные сочетания угла и спектра подсветки дают радикально несхожие изображения одной надписи, и эта чувствительность образует ещё один источник доменного сдвига, действующий на уровне процедуры съёмки.

\subsection{Инфраструктура цифровых рунических корпусов}
\label{sec:corpora}

\subsubsection{Samnordisk runtextdatabas (Rundata)}

База Samnordisk runtextdatabas (далее~-- Rundata), известная также как Scandinavian Runic-text Database (SRD), является стандартным транскрипционным корпусом скандинавских рунических надписей. Проект возник в Уппсальском университете на рубеже 1980--1990-х годов: первая версия собрана в 1987~г.\ на основе локальной базы шведских надписей, а финансирование для расширения до общенордического охвата получено в 1992~г.~\cite{uu2024rundata}. Корпус последовательно эволюционировал (Rundata~1.0, 2001; 2.5, 2007; 3.0, 2008; 3.1, январь~2018), и с декабря~2020~г.\ основная пользовательская поддержка перенесена на веб-платформу Runor под управлением Шведского управления по охране культурного наследия (Riksantikvarieämbetet).

Ограничение Rundata для задач компьютерного зрения состоит в том, что подавляющее большинство записей содержит только текстовую информацию~-- транслитерацию, нормализацию, перевод, библиографические ссылки и другую метаинформацию~-- и не имеет привязанных машиночитаемых изображений. Изображения отдельных надписей частично доступны через интерфейс Runor, где Riksantikvarieämbetet интегрирует фотодокументацию из собственных и партнёрских архивов~\cite{riksant2024}, однако ни покрытие, ни разрешение, ни условия съёмки этих изображений не унифицированы, и формальной связки «текстовая запись\,$\leftrightarrow$\,массив изображений с фиксированными метаданными» в базе не предусмотрено. Это означает, что прямое использование Rundata как обучающего датасета CV-модели невозможно: база предоставляет богатую целевую разметку (ground truth по транслитерации, нормализации и атрибуции), но без предварительной и основательной ручной обработки данные не применимы.  Этим же ограничением объясняется, почему в рецензируемой литературе не выявлено ни одной ML-публикации, использующей Rundata как тренировочный материал, несмотря на то что сообщество явно признает такую возможность как нерешённую инфраструктурную задачу~\cite{williams2022,magin2024}.

\subsubsection{Альтернативные базы и корпусные ресурсы}

Вторую цифровую инфраструктуру предоставляет проект RuneS («Runic Writing in the Germanic Languages») Союза академий наук Германии (Akademie der Wissenschaften zu Göttingen) действующий с 2010~г.\ и охватывающий руническое письмо примерно со 100 по 1500~г.\ н.\,э.~\cite{duwelmarold2011,runesdb}. В отличие от Rundata, фокусированной на скандинавском ареале, RuneS систематически включает южногерманский корпус, англо-фризский футорк и руническую рукописную традицию (\textit{runica manuscripta}). База содержит отдельный модуль аллографов и графотипических вариантов с описанием формы знака на уровне отдельных штрихов, чего в Rundata нет. Технически RuneS реализован как веб-интерфейс с возможностью браузерного поиска \texttt{runesdb}. Лицензионный режим напрямую допускает использование данных в обучении ML-моделей при атрибуции и сохранении лицензии~\cite{runesdb}. При этом RuneS, как и Rundata, не предоставляет в открытом доступе унифицированных изображений, пригодных для обучения CV-моделей.

Nationalmuseet Danmark поддерживает базу Danske Runeindskrifter (свыше 900 датских надписей с фотодокументацией)~\cite{nationalmuseet2024}; Kulturhistorisk museum в Осло поддерживает Runic Archives (свыше 1600 норвежских надписей)~\cite{khmoslo2024}; The British Museum предоставляет открытый коллекционный API с фотографиями отдельных англосаксонских артефактов; Europeana функционирует как наднациональный агрегатор кросс-институционального поиска~\cite{europeana2024}. Качество и формат изображений в этих репозиториях различны: разрешение варьирует от низкокачественных архивных сканов начала 2000-х до 6000-пиксельных современных фотографий; условия освещения не унифицированы, а оптимальная для эпиграфики raking light-съёмка применяется крайне редко.

\subsubsection{Архитектурные ограничения для машинного обучения}
\label{sec:ml-architecture-limits}

Перечисленные ресурсы создают для машинного обучения четыре связанных ограничения.

Первое~-- разрыв между разметкой и сигналом, описанный выше: целевая транслитерация существует в Rundata и RuneS, но входные изображения приходится извлекать и предобрабатывать отдельно, что делает связывание по \texttt{signum} обязательным предварительным этапом.

Второе~-- неравномерность распределения по системам и периодам. Распределение записей радикально несбалансировано: на Младший футарк приходится подавляющее большинство записей, Старший футарк представлен примерно на порядок меньшим числом надписей; Англосаксонский футорк~-- менее чем двумястами~\cite{page1999,riksant2024}. Это распределение отражает археологические факторы сохранности, и его нельзя выровнять добавлением новых записей, только посредством обогащения выборки синтетическими данными. Class imbalance имеет двухуровневую структуру~-- на уровне систем и на уровне отдельных графем внутри системы~-- и, как показано в типологии доменного сдвига, напрямую транслируется в диахронический сдвиг между сегментами выборки. Обучение на полной выборке без коррекции сводит baseline-предсказание к мажоритарному классу Младшего футарка, что требует взвешенной функции потерь, стратифицированной выборки либо раздельных моделей с system-aware classifier head.

Третье~-- несовместимость нотаций между базами. Одна и та же надпись в Rundata, RuneS и национальных каталожных изданиях может иметь различную транслитерацию~-- не из-за ошибок, а из-за различий редакционных конвенций. Williams, Bianchi и Zimmermann~\cite{williams2022} документируют расхождения и показывают, что для отдельных надписей четыре основных источника могут давать четыре несовпадающие транслитерации, отличающиеся не только в спорных позициях, но и в общем числе графем. Отсюда требование явного выбора единого референсного стандарта  и ведения «лога расхождений»~-- сохранения альтернативных транслитераций как дополнительных полей.

Четвёртое~-- ограничения Unicode. Блок Runic (U+16A0--U+16FF), введённый Unicode~3.0 (1999) и расширенный восемью знаками в Unicode~7.0 (2014), содержит 89~кодпойнтов, унифицирующих графемы трёх рунических систем, средневековые скандинавские руны и символы рунических календарей~\cite{unicodeRunic}. Принципиальное ограничение, разбираемое у Magin и Smith~\cite{magin2024}, в том, что кодпойнт-уровень не различает аллографы: один кодпойнт покрывает все графические варианты знака~-- закруглённые, угловатые, зеркальные, с разной длиной плеч. Текущая реализация блока несовместима ни с задачей структурированного поиска по аллографам, ни с задачей классификации, различающей аллографы как отдельные классы; авторы констатируют, что спустя почти 25~лет после введения блока рунологи в массе своей им не пользуются~\cite{magin2024}. Для машинного обучения это означает, что Unicode-кодпойнт не может служить меткой классификатора аллографов. В экспериментальной части исследования была предпринята попытка дообучить OCR на транслитерированной версии текста, то есть без использования рун из блока Runic. 

\subsubsection{Стратегия интеграции в ML-конвейер}

\hl{С учётом этих ограничений принимается следующая схема сборки датасета. Из Rundata загружается полный экспорт текстовых записей в формате клиента~3.1, парсимый с восстановлением полей \texttt{signum}, локации, датировки, материала, рунической системы, транслитерации и нормализации; параллельно через пакет \texttt{runesdb} загружаются записи RuneS для тех \texttt{signum}, что присутствуют в обеих базах, ради альтернативных транслитераций и аллографической разметки. Далее текстовые записи связываются с изображениями: для шведских надписей~-- через K-samsök\,/\,Söktjänsten Runor по сопоставлению \texttt{signum}; для датских~-- через Danske Runeindskrifter; для норвежских~-- через KHM Runic Archives; для англосаксонских~-- через коллекции The British Museum. Связка проходит ручную валидацию: protocol-файл фиксирует для каждой пары «\texttt{signum}\,$\leftrightarrow$\,изображение» уровень уверенности (verified\,/\,probable\,/\,uncertain) и причину отбраковки.

Каждой валидированной паре присваивается унифицированный внутренний идентификатор формата \texttt{RUN-XXXXXX} и UUID4 для использования в инструментах ML-конвейера (Weights~\&~Biases, MLflow, HuggingFace Datasets). Введение собственного идентификатора мотивируется тремя соображениями. Во-первых, \texttt{signum} Rundata не является глобально уникальным: одна физическая надпись может фигурировать под несколькими \texttt{signum} при пересмотре атрибуции, и собственный идентификатор обеспечивает стабильную точку привязки. Во-вторых, одна запись может соответствовать нескольким изображениям (разное освещение, общий план и крупные планы фрагментов), а задача требует первичной привязки к изображению, поэтому первичной единицей корпуса является пара «\texttt{signum}\,+\,изображение». В-третьих, в корпус включаются англосаксонские и континентальные надписи без \texttt{signum} в Rundata, и единая схема позволяет рассматривать их наравне со скандинавскими.}

Таблица~\ref{tab:digital-resources} систематизирует характеристики основных ресурсов.

\begin{table}[htbp]
\centering
\caption{Сравнительная характеристика рунических цифровых ресурсов}
\label{tab:digital-resources}
\footnotesize
\begin{tabular}{p{2.0cm} p{2.3cm} p{2.4cm} p{2.4cm} p{3.2cm}}
\hline
\textbf{База} & \textbf{Институция} & \textbf{Объём / изображения} & \textbf{Лицензия} & \textbf{Применимость для ML} \\
\hline
Rundata (SRD) & Uppsala Univ. & $\sim$7000 записей; изображений в базе нет & ODbL\,+\,DbCL & целевая разметка без входного сигнала; нужно связывание \\
RuneS & Göttingen, Kiel, Eichstätt & Старший, англо-фризский, Младший футарк; без изображений & CC-BY-SA 4.0 & детальная палеографическая и аллографическая разметка \\
Söktjänsten Runor & Riks\-antikvarie\-ämbetet & $\sim$7000 агрегир.; изображения фрагментарны & CC0 (метаданные); CC-BY (изобр.) & основной доступный источник изображений с привязкой к \texttt{signum} \\
Historiska / K-samsök & SHM, RAÄ & $\sim$5,7 млн изображений всех типов & CC0; CC-BY(-SA) & основной источник изображений шведских артефактов \\
Europeana & Europeana Foundation & агрегатор провайдеров & по провайдеру & удобен для кросс-страновых выборок \\
\hline
\end{tabular}
\end{table}

Описанная схема образует основу для построения корпуса в Главе~3. Прежде чем перейти к разметке, необходимо решить, какой транслитерационный стандарт принимается за референс, поскольку выбор между Rundata, RuneS и национальными конвенциями содержательно меняет целевую метку модели и численную интерпретацию метрик CER. Этому посвящён следующий раздел.

\subsection{Конвенции транслитерации и формализация эпиграфической неопределённости}
\label{sec:transliteration}

Прежде чем обучать модель на рунических данных, необходимо решить вопрос, 
который не возникает в стандартных OCR-задачах: что именно 
считать правильным ответом. В рунологии необходимо рассматривать его сразу на двух уровнях. На первом -- выбор транслитерационного 
стандарта: параллельно существует несколько конкурирующих конвенций, 
каждая из которых по-своему кодирует знаки Младшего футарка 
и систему обозначений неуверенного прочтения. На втором -- сама надпись: 
даже зафиксировав единый стандарт, мы сталкиваемся с тем, что значительная 
доля надписей допускает несколько равноправных прочтений, между которыми 
консенсус рунологов не достигнут.

Для машинного обучения оба фактора имеют прямые методологические последствия. 
Если целевая метка $y$, относительно которой измеряется отклонение предсказания, 
оказывается недоопределённой, то интерпретируемость стандартных метрик -- CER, 
sequence accuracy, F1 по графемам -- теряется: численное значение ошибки 
зависит от того, какое именно из допустимых прочтений выбрано в качестве 
эталона. В разделе формализуются оба источника неопределённости и описываются конвенции, принимаемые в настоящей работе.

\subsubsection{Множественность стандартов транслитерации}

Транслитерация (конверсия «руна\,$\to$\,латинская буква или диграф») принципиально отличается от транскрипции, то есть нормализации к стандартному древнескандинавскому или древнеанглийскому~\cite{spurkland2005,barnes2012}. Она формально однозначна по построению: каждой графеме приписывается каноническое латинское соответствие безотносительно реально передававшегося звука. Однако, в рунологии устоялась не одна традиция транслитерации.

Конвенция Rundata опирается на принципы, восходящие к работам Лены Петерсон~\cite{jesch2013}. Для Младшего футарка она принимает «уплощённую» нотацию: \textit{kaun} систематически транслитерируется как \textbf{k}, хотя передаёт \textipa{/k/}, \textipa{/g/} и \textipa{/ŋ/}; \textit{týr}~-- как \textbf{t}, передавая \textipa{/t/} и \textipa{/d/}; \textit{ár}~-- как \textbf{a}, передавая \textipa{/a/}, \textipa{/æ/}, \textipa{/e/} и назализованный \textipa{/ã/}~\cite{spurkland2005,duwel2008}. Нотация повреждённых прочтений кодифицирована так: квадратные скобки~-- восстановленная часть; курсив или подчёркивание~-- неуверенные знаки; знак вопроса~-- альтернативное прочтение; многоточие~-- лакуна; круглые скобки~-- редакторские дополнения~\cite{jesch2013}. Эта система, удобная для филологического чтения, при буквальном переносе в обучающую разметку создаёт сложности: специальные символы интерпретируются токенайзером как обычные графемы и портят словарь меток.

Классическая система Краузе--Янкуна (KJ), опубликованная в \textit{Die Runeninschriften im älteren Futhark}~\cite{krause1966}, остаётся стандартом транслитерации Старшего футарка и отличается от Rundata большей детализацией: последовательнее разделяет аллографы, использует диакритические знаки для особых рун и систематически применяет звёздочку для реконструированных форм~\cite{krause1966,antonsen1975}. Поскольку для записей Старшего футарка Rundata в значительной мере воспроизводит KJ-конвенции, а собственная конвенция Rundata оптимизирована под Младший футарк, корпус демонстрирует внутреннюю стилистическую неоднородность транслитераций между сегментами.

Принципы RuneS опубликованы как часть документации базы~\cite{runesdb}. Они агрегируют разные источники с сохранением происхождения: транслитерации Старшего футарка заимствуются из Runenprojekt Kiel, датские надписи Младшего футарка~-- из Danske Runeindskrifter, прочие скандинавские~-- из Rundata~2.5, англо-фризские~-- собственные принципы RuneS. RuneS не унифицирует эти конвенции в единый стандарт, что для машинного обучения означает дополнительный источник неоднородности целевых меток. Национальные печатные стандарты служили референсами для Rundata, но каждый использует собственные конвенции передачи специальных знаков и обозначения лакун~\cite{williams2022}.

\subsubsection{Эпиграфическая неопределённость как проблема ground truth}

Даже после фиксации единого стандарта значительная доля надписей сохраняет неустранимую неопределённость прочтения, имеющую несколько источников. Во-первых, физическая повреждённость: эрозия, сколы и трещины приводят к тому, что отдельные графемы видны лишь частично, и полная и однозначная трактовка знака невозможна. Во-вторых, графематическая полифония Младшего футарка, где одна графема соответствует нескольким фонемам и выбор возможен лишь на уровне контекстной интерпретации. В-третьих, аллографическая интерпретация: на пограничных случаях длинноветвистая и коротковетвистая формы допускают двойное отнесение, и разные рунологи классифицируют один физический рез как реализацию разных графем. В-четвёртых, чтение направления и порядка: в спирально или дуговидно расположенных надписях направление устанавливается лишь с привлечением иконографического и формульного анализа~\cite{kallstrom2007,williams2010}.

Эпиграфическая неопределённость не может быть устранена ни более качественным изображением, ни лучшей моделью~-- она встроена в природу источника. Это означает, что классическое понятие ground truth, являющееся базовым для машинного обучения, применительно к рунической OCR оказывается недоопределённым.

\subsubsection{Формализация неопределённости для машинного обучения}

Первый возможный подход~-- жёсткая фильтрация: из обучающей выборки исключаются все надписи с маркерами неуверенного прочтения. Подход прост и даёт корректно интерпретируемый CER на оставшейся выборке, но сильно сокращает объём данных и систематически смещает распределение в сторону лучше сохранившихся камней; обученная на этом подмножестве модель может показывать заниженный CER на in-distribution-тесте и резко более высокий~-- на повреждённых надписях, к работе с которыми и стремится практика.

Второй подход~-- мягкая разметка (soft labeling): метка превращается из единственной строки в распределение вероятностей по альтернативным прочтениям. На уровне отдельного знака мягкая метка задаётся как распределение $p_i(c)$ по множеству допустимых графем:
\begin{equation}
p_i(c) = \begin{cases}
1, & c = c^*_i \text{~-- уверенное прочтение} \\
\dfrac{1}{|\mathcal{A}_i|}, & c \in \mathcal{A}_i \text{~-- множество альтернатив} \\
0, & \text{иначе}
\end{cases}
\label{eq:softlabel}
\end{equation}
где $\mathcal{A}_i$~-- экспертно-определённое множество альтернативных графем в позиции $i$ для неуверенно прочитываемого знака. Кросс-энтропийная функция потерь обобщается на soft labels:
\begin{equation}
\mathcal{L}_{\text{soft}}(\hat{p}, p) = -\sum_{i=1}^{T} \sum_{c \in \mathcal{C}} p_i(c) \log \hat{p}_i(c),
\label{eq:softloss}
\end{equation}
где $\hat{p}_i(c)$~-- предсказанная вероятность графемы $c$ в позиции $i$, $\mathcal{C}$~-- полный графемный инвентарь, $T$-- длина последовательности. Такая функция потерь не штрафует модель за предсказание любого из допустимых альтернативных прочтений, что соответствует признанию обоих вариантов экспертом.

Соответственно переопределяется и метрика. Стандартный CER рассчитывается как
\begin{equation}
\text{CER} = \frac{S + D + I}{N},
\label{eq:cer}
\end{equation}
где $S$, $D$, $I$~-- числа замен, удалений и вставок в редакционном расстоянии между предсказанной и эталонной строками, $N$~-- длина эталона. В случае мягкой разметки CER заменяется на ambiguity-aware CER (A-CER), в котором подстановка между двумя графемами, признанными альтернативами в данной позиции, считается не ошибкой, а корректным предсказанием:
\begin{equation}
\text{A-CER} = \frac{S^* + D + I}{N}, \qquad S^* = \sum_{i} \mathbb{1}\bigl[\hat{c}_i \notin \mathcal{A}_i\bigr],
\label{eq:acer}
\end{equation}
где $\hat{c}_i$~-- предсказание модели в позиции $i$, а индикатор $\mathbb{1}[\cdot]$ возвращает единицу только тогда, когда предсказанная графема не входит в множество допустимых альтернатив. Эта метрика, идейно восходящая к обобщениям CER в литературе по HTR исторических документов, применительно к рунической эпиграфике формализуется здесь явно и используется параллельно со стандартным CER для разделения двух типов ошибок: подлинных ошибок распознавания и допустимых альтернатив.

Изложенный аппарат (1)–(2),(4) представляет формализацию, которую работа вносит как методологический задел. В текущей базовой постановке сознательно принята жёсткая разметка и стандартный CER (4), поскольку цель ~-- установить воспроизводимый baseline на однозначных прочтениях, тогда как A-CER требует экспертно-размеченного множества альтернатив​, формирование которого вынесено в перспективы,»

\subsubsection{Принятые в работе соглашения}

Основным стандартом транслитерации выбрана конвенция Rundata. Важно, что транслитерация в этой конвенции выступает в работе не вспомогательным представлением, а целевым выходом системы распознавания: модель предсказывает латинскую строку Rundata (например, \texttt{fuþark}), а не последовательность рунических кодовых позиций Unicode. Такой выбор устраняет необходимость расширения словаря токенизатора рунами блока U+16A0--U+16FF и позволяет переиспользовать предобученные субсловные словари без модификации, что критично в условиях малого объёма данных. 

Предобработка исходных транслитераций в обучающий корпус строится как нормализация к чистому графемно-транслитерационному ядру конвенции Rundata. Вся эпиграфическая разметка~-- разделители слов, скобочные обозначения лакун и неуверенных прочтений, маркеры лигатур (\texttt{=}) и альтернативных чтений (\texttt{/}), пометы строк и секций~-- трактуется как граница токена; сохраняются лишь токены, целиком состоящие из символов допустимого транслитерационного алфавита (33 знака, включая \texttt{þ}, \texttt{ð}, \texttt{æ}, \texttt{ø}, \texttt{ą}, носовые и \texttt{R}/\texttt{ʀ}), при этом заглавная \texttt{A} приводится к строчной, а \texttt{R} и \texttt{ʀ} сохраняются как различные графемы. Для генерации синтетических изображений очищенная транслитерация детерминированно отображается обратно в руны полной обратимой таблицей соответствий, целевой инвентарь которой составляют младшие руны long-branch как ядро, пунктированные руны для дополнительных фонем (\texttt{e}, \texttt{g}, \texttt{d}, \texttt{p}, \texttt{y}, \texttt{o}, \texttt{c}, \texttt{z}, \texttt{v}) и отдельные старшие/англосаксонские графемы для символов без младше-футаркового эквивалента (\texttt{w}, \texttt{j}, \texttt{ð}, \texttt{ï}, \texttt{ŋ}, \texttt{æ}, \texttt{ø}). Обратимость отображения обеспечивает строгое соответствие между изображением и меткой по построению, что исключает рассогласование разметки как источник шума в синтетических данных.

Из принятой процедуры следует явное разделение режимов учёта эпиграфической неопределённости. Синтетический корпус формируется из уверенных чистых прочтений, неопределённость и лакуны в него намеренно не вносятся, а его метки точны, что позволяет оценивать распознавание на нём стандартной метрикой. 

Принятая совокупность соглашений~-- конвенция Rundata как стандарт транслитерации и одновременно цель распознавания, нормализация реальных транслитераций к чистому алфавиту, детерминированное обратимое отображение для синтеза и разделение режимов неопределённости между синтетическими и реальными данными~-- образует формальный аппарат, относительно которого далее описываются схема разметки корпуса, процедура его сборки и протокол оценки.


\section{Эмпирическая база и формирование корпуса}
\label{sec:data}
 
Целевая разметка скандинавских надписей существует, однако привязанных к ней унифицированных машиночитаемых изображений нет, а связывание текстовых записей с фотографиями по \texttt{signum} требует ручной валидации каждой пары. Из этого ограничения следует, что собрать реальный изобразительно-текстовый корпус объёмом, достаточным для обучения нейросетевого распознавателя «с нуля», что в рамках классического OCR требует десятки, а иногда и сотни тысяч верифицированных пар, в рамках настоящей работы невозможно. 
 
\subsection{Стратегия формирования данных: обучение на синтетике, оценка на реальном золотом стандарте}
\label{sec:data-strategy}
 
Эмпирическая база работы строится на принципиальном разделении источника обучающего сигнала и источника оценочного сигнала. Обучающая выборка формируется полностью синтетически: изображения рунических надписей генерируются управляемым конвейером, а их транслитерационные метки точны по построению благодаря обратимому отображению. Оценочная выборка, напротив, состоит исключительно из реальных фотографий надписей с выверенной транслитерацией и используется только для валидации и тестирования, но не для обучения. 

Наиболее близкая по постановке работа Derin и Uçar~\cite{derin2024} обучает Vision Transformer на полностью синтетическом корпусе древнетюркских руноподобных знаков, полученном параметрической рандомизацией и физически-основанным рендерингом, и применяет его в zero-shot-режиме к реальным фотографиям камня Кюль-Тегина с высокой точностью без единого реального обучающего примера. Тем самым связка «обучение на синтетике~-- оценка на реальных данных» для руноподобной эпиграфики является не обходным приёмом, а воспроизведением подтверждённой исследовательской парадигмы. Её дополняют два смежных результата: Vogler и др.~\cite{vogler2022} демонстрируют, что для адаптации предобученного энкодера к новой письменной традиции достаточно нескольких десятков размеченных строк, что оправдывает минимизацию объёма реальных данных; Dencker и др.~\cite{dencker2020} используют латинские транслитерации как выравнивающий сигнал для разметки знаков, что обосновывает выбор реальных транслитераций Rundata в качестве лексической основы синтеза.
 
В условиях, когда число доступных верифицированных реальных пар измеряется десятками, каждая такая пара обладает наибольшей ценностью именно как несмещённая точка оценки обобщающей способности, а не как один из многих обучающих примеров. Включение реальных данных в обучение одновременно загрязнило бы тестовую оценку утечкой распределения и одновременно сократило бы и без того малую тестовую выборку. Удержание всего реального множества вне обучения позволяет измерить величину, составляющую центральный эмпирический вопрос работы,~-- разрыв между синтетическим и реальным доменами (domain shift) для рунического письма. 

Помимо этого, синтетический корпус формируется из уверенных чистых прочтений с точными метками, и качество на нём измеряется стандартным CER. Поскольку реальный верифицированный набор данных в рамках принятой методологии также приводится к однозначной строковой репрезентации, его валидация осуществляется по аналогичной схеме с использованием классического посимвольного расстояния редактирования. Это обеспечивает полную метрическую унификацию между синтетическим и реальным доменами, позволяя напрямую оценивать величину падения точности при переносе модели.
 
Такой подход позволяет корректно интерпретировать целевой порог CER. Порог в 20\,\% понимается не как безусловный результат, а как априорно фиксируемая граница практической применимости: распознавание с долей посимвольных ошибок ниже этого уровня сохраняет полезность для индексирования, поиска по корпусу и последующей экспертной корректуры, тогда как обучение модели исключительно на синтетических данных оставляет открытым эмпирический вопрос о том, насколько близко к этой границе удаётся подойти при переносе на реальные носители. Порог 20\,\% принимается как рабочая инженерная граница, ниже которой выход пригоден для индексации и экспертной пост-корректуры. Он не претендует на универсальность и служит единой точкой отсчёта для интерпретации результатов. Соответственно, вклад экспериментальной части состоит не в гарантированном достижении порога, а в установлении первой воспроизводимой базовой линии (baseline) распознавания для рунического материала и в количественной характеризации синтетически-реального разрыва как величины, определяющей дальнейшие направления его сокращения.
 
\subsection{Золотой стандарт (gold set)}
\label{sec:gold}
 
Золотым стандартом (gold set) в работе называется компактное множество реальных фотографий рунических надписей с экспертно выверенной транслитерацией, используемое исключительно для валидации и тестирования. Его объём сознательно ограничен порядком нескольких десятков-сотен надписей: как показано в подразделе, бóльшая выборка верифицированных пар «\texttt{signum}\,$\leftrightarrow$\,изображение» недостижима без непропорциональных затрат ручной разметки, тогда как для несмещённой оценки переносимости модели и для типологического анализа ошибок выборки такого порядка достаточно при условии корректного учёта статистической погрешности. 
 
Транслитерационные метки gold set полностью нормализуются к чистокровному графемно-транслитерационному ядру конвенции Rundata. Вся вспомогательная эпиграфическая разметка (индикаторы лакун, разделители, повреждения носителя) отсекается на этапе предобработки, формируя жесткую целевую последовательность (hard label), состоящую исключительно из знаков допустимого унифицированного алфавита.

В отличие от подходов, допускающих вероятностное сглаживание меток (soft labeling), в данной конфигурации пайплайна для каждой знаковой позиции на реальном носителе выбирается строго один целевой графемный класс, соответствующий наиболее аргументированному консенсусному прочтению в базе данных Rundata. Документированные расхождения транслитераций между эпиграфическими школами выносятся за рамки графематического ядра модели; они фиксируются в изолированном «логе расхождений» и привлекаются исключительно на этапе качественного интертекстуального анализа ошибок, не влияя на расчет целевой функции потерь и метрики CER.
 
Разметка золотого стандарта организована иерархически на трёх уровнях, что обеспечивает совместимость с задачами компьютерного зрения и с последующим лингвистическим слоем и допускает повторное использование корпуса. Верхний уровень~-- надписи (\textit{inscription})~-- фиксирует метаданные памятника: \texttt{signum} и внутренний идентификатор, географическую локализацию, датировку с указанием её точности, тип материала и состояние сохранности, руническую систему, библиографический источник и параметры съёмки (тип освещения, разрешение). Средний уровень~-- строки (\textit{line})~-- задаёт порядок и направление чтения и область строки на изображении. Нижний уровень~-- графемы (\textit{glyph})~-- содержит транслитерационную графему, помету аллографа (long-branch\,/\,short-twig\,/\,staveless), степень уверенности прочтения и, при наличии, множество альтернатив $\mathcal{A}_i$.
 
\begin{table}[htbp]
\centering
\caption{Иерархическая схема аннотации золотого стандарта}
\label{tab:annotation-schema}
\small
\begin{tabular}{p{2.6cm} p{3.7cm} p{6.4cm}}
\hline
\textbf{Уровень} & \textbf{Единица} & \textbf{Поля разметки} \\
\hline
Надпись & памятник (\textit{inscription}) & \texttt{signum}, внутренний ID (\texttt{RUN-XXXXXX}), локализация, датировка и её точность, материал, сохранность, рунич.\ система, источник, параметры съёмки (освещение, разрешение) \\
Строка & строка / сегмент (\textit{line}) & порядок и направление чтения, область строки на изображении \\
Графема & знак (\textit{glyph}) & транслитерационная графема, аллограф, степень уверенности, множество альтернатив $\mathcal{A}_i$ \\
\hline
\end{tabular}
\end{table}
 
Отбор изображений в золотой стандарт ведётся не пропорционально распределению корпуса, а с ориентацией на покрытие источников гетерогенности, порождающих доменный сдвиг: рунической системы (Старший, Младший, Англосаксонский футарк), типа носителя (камень, металл, кость) и условной градации сохранности. При этом ограниченность представленности редких систем явно фиксируется как предел покрытия золотого стандарта, а не как свойство модели.
 
Назначение gold set остаётся строго оценочным: он образует тестовое множество и не участвует ни в обучении, ни в подборе гиперпараметров, для которых используется отложенная синтетическая подвыборка; при необходимости честной ранней остановки из него может выделяться небольшая девелопмент-доля, не пересекающаяся с тестовой. Ограниченный объём множества прямо влияет на интерпретацию результатов: оценка CER на выборке такого порядка сопровождается широкими доверительными интервалами, в связи с чем в экспериментальной части она приводится с бутстрэп-оценкой погрешности и трактуется как индикативная, а не как точечная характеристика качества.


\subsection{Синтетический корпус}
\label{sec:synth}
 
Если золотой стандарт обеспечивает оценочный сигнал, то весь обучающий сигнал работа черпает из синтетического корпуса. Корпус должен одновременно удовлетворять двум требованиям, которые описаны в данной главе. Во-первых, лексическая достоверность: изображённые знаковые последовательности обязаны быть реальными руническими словами, а не случайными наборами графем, иначе распознаватель обучится на несуществующем распределении языка. Во-вторых, визуальная достоверность деградированного носителя: структурный, а не аддитивный характер шума требует, чтобы синтетическое изображение воспроизводило именно те признаки локального контраста и текстуры, на которых распознаватель учится отделять рез от фона. Первое требование удовлетворяется заимствованием лексики из реальных транслитерационных баз; второе~-- управляемым диффузионным синтезом с обусловливанием по контуру.
 
\subsubsection{Лексическая основа: извлечение слов из реальных баз}
\label{sec:synth-lexicon}
 
Лексика синтетического корпуса извлекается из текстовых экспортов реальных рунических баз (записи, восходящие к Rundata, и сводка Gamla runor Кристера Хампа), что прямо реализует логику Dencker и др.~\cite{dencker2020}: транслитерация выступает выравнивающим сигналом, связывающим знак и его латинское соответствие. Колонка транслитераций токенизируется на отдельные слова по границам, заданным символами вне допустимого алфавита; сохраняются токены длиной от трёх до двадцати пяти знаков, целиком состоящие из символов чистого графемно-транслитерационного ядра конвенции Rundata. Каждое слово отображается в руническую строку полной обратимой таблицей соответствий, после чего тождество \textit{транслитерация}~$\to$~\textit{руны}~$\to$~\textit{транслитерация} проверяется на всём корпусе: совпадение по построению гарантирует, что метка-строка точна.
 
Чтобы отсечь разовый шум транскрипции, в корпус включаются лишь слова, встретившиеся не менее двух раз; их частоты задают распределение, из которого ведётся выборка с температурой~0,7, так что частотные слова представлены чаще, приближая статистику обучающей выборки к корпусной. Из слов собираются фразы длиной от одного до четырёх слов, разделяемые рунической точкой-разделителем (\texttt{U+16EC}, ᛬), что воспроизводит характерное для каменных надписей разделение слов и снабжает распознаватель примерами межсловных границ.
 
\subsubsection{Рендеринг рунической строки}
\label{sec:synth-render}
 
Полученная руническая строка отрисовывается шрифтом Noto Sans Runic чёрным по белому на холсте $1024\times1024$ с автоматическим подбором кегля: размер шрифта понижается до тех пор, пока строка не умещается в рабочую область, а при невозможности уместить её в одну строку выполняется перенос на несколько строк. Небольшое горизонтальное дрожание положения вносит позиционную вариативность. 

Здесь же фиксируется первое ограничение покрытия синтетического корпуса. Шрифт реализует кодовые позиции рунического блока Unicode, а кодпойнт не различает аллографы: одна позиция покрывает все графические варианты знака. Следовательно, синтетические глифы представляют единственное аллографическое семейство на кодпойнт и не воспроизводят полную вариативность long-branch\,/\,short-twig\,/\,staveless. Это сужение~-- самостоятельный источник синтетически-реального разрыва, и его расширение (несколько рунических шрифтов, аллографические формы RuneS) обсуждается как направление развития.
 
\subsubsection{Выбор парадигмы синтеза и модальности обусловливания}
\label{sec:synth-paradigm}
 
Центральное проектное решение синтетического конвейера~-- способ превращения отрисованной строки в фотореалистичное изображение каменной надписи. Пространство возможных парадигм сводится к пяти, и их сопоставление по трём критериям~-- контроль над целевой меткой, реалистичность врезки знака и потребность в реальных данных целевого письма~-- сведено в таблицу~\ref{tab:synth-paradigms}.
 
\begin{table}[htbp]
\centering
\caption{Сопоставление парадигм синтеза рунических изображений}
\label{tab:synth-paradigms}
\small
\begin{tabular}{p{3.3cm} p{2.5cm} p{2.3cm} p{2.0cm} p{2.6cm}}
\hline
\textbf{Подход} & \textbf{Контроль метки} & \textbf{Реалистичность врезки} & \textbf{Нужны реальные данные письма} & \textbf{Итог} \\
\hline
Text-to-image без управления & нет & высокая & нет & непригоден: прочтение не задаётся \\
GAN-синтез письма & да & средняя & да (корпус целевого письма) & непригоден: требует отсутствующего корпуса \\
Параметрический 3D-рендеринг & точный & высокая & нет & эталон, но высокая инженерная стоимость \\
Диффузия + inpainting (маска рун) & точный & средняя & нет & принят как геометрически точный режим \\
Диффузия + ControlNet Canny & приближённый & высокая & нет & принят как основной режим \\
\hline
\end{tabular}
\end{table}
 
Два подхода исключаются по построению. Свободный синтез text-to-image выдаёт фотореалистичную «руноподобную» поверхность, но не позволяет задать, какие именно руны будут изображены, а значит, не порождает известной метки и непригоден для супервизорного обучения. GAN-синтез письма (линия ScrabbleGAN~\cite{fogel2020scrabblegan}) допускает обусловливание текстом, однако обучается на реальном корпусе целевого письма, отсутствие которого и составляет исходную проблему работы; для рунического материала с десятками реальных образцов такая модель неустойчива. Параметрический трёхмерный рендеринг знака, врезанного в смоделированную каменную поверхность с рандомизацией освещения и выветривания, даёт одновременно точную геометрию, реалистичную врезку и точную метку и служит эталоном в работе Derin и Uçar~\cite{derin2024}; его недостаток~-- высокая инженерная стоимость построения трёхмерных ассетов, материалов и шейдеров выветривания, плохо совместимая с одиночной разработкой в условиях вычислительных ограничений.
 
Остаются два режима на основе предобученной латентной диффузии~\cite{rombach2022ldm}, не требующие дообучения под рунический скрипт и потому естественные для малоресурсной задачи; они различаются по оси «точность геометрии~-- реалистичность врезки». В режиме inpainting отрисованные руны маскируются вместе с буфером дилатации, и генерация физически не затрагивает зону знака: геометрия штрихов сохраняется попиксельно, а метка точна, однако руны остаются плоскими отрисованными штрихами, скомпонованными с синтезированным фоном, и читаются скорее как «нанесённые поверх», нежели как врезанные, без собственных теней и рельефа. В режиме ControlNet~\cite{zhang2023controlnet} всё изображение, включая зону знака, синтезируется заново под управлением контурной карты, и руны выходят согласованно врезанными в камень с единым освещением и выветриванием. Однако в таком режиме возможны артефакты на изображениях в виде неправильного начертания одной или нескольких рун. Такой метод требует ручной выборочной валидации. Также замечено, что появление нежелательных артефактов присуще изображениям с короткими текстами (3-4 символа). Работа использует оба режима: ControlNet~-- как основной, обеспечивающий фотореалистичную врезку, inpainting~-- как геометрически точный комплемент, прямое сопоставление которых выносится в абляционное исследование экспериментальной части.
 
Внутри парадигмы ControlNet остаётся выбор модальности обусловливания. ControlNet допускает управление по различным картам признаков, и их пригодность для тонких штрихов руны существенно различается; сопоставление приведено в таблице.
 
\begin{table}[htbp]
\centering
\caption{Сопоставление модальностей обусловливания ControlNet для рунической геометрии}
\label{tab:controlnet-modalities}
\small
\begin{tabular}{p{2.6cm} p{3.4cm} p{3.0cm} p{4.0cm}}
\hline
\textbf{Модальность} & \textbf{Характер сигнала} & \textbf{Пригодность для штриха} & \textbf{Причина отклонения / выбора} \\
\hline
Canny & жёсткие бинарные контуры & высокая: чёткие границы каждого штриха & \textbf{выбрана}: точная топология при управляемой жёсткости \\
HED / soft-edge & мягкие размытые края & ниже: размывает границы тонких штрихов & больше свободы модели $\to$ сильнее дрейф геометрии \\
Scribble & эскизный, очень свободный & низкая & избыточная свобода интерпретации формы \\
Lineart & тонкая чистая линия (иллюстрация) & средняя, но доменно несовместима & прайор иллюстративной графики, мягче Canny \\
Depth / normal & карта глубины / нормалей (3D) & неинформативна для плоского рендера & у плоских рун нет глубины; требует синтеза рельефа (см.\ \ref{sec:synth-fidelity}) \\
Segmentation & семантические области & неприменима к штрихам & оперирует классами регионов, не линиями \\
MLSD & только прямые отрезки & частичная: теряет кривые & искажает инвентарь (кривые формы, разделитель) \\
\hline
\end{tabular}
\end{table}
 
Выбор Canny мотивирован тем, что руническая графема определяется прежде всего топологией своих штрихов~-- ствола и ветвей,~-- которую бинарная карта рёбер фиксирует чище всего: на контрастном рендере детектор Canny выделяет однозначные границы каждого штриха, задавая жёсткий структурный каркас. Остальные модальности уступают по характеру сигнала. Мягкие края HED размывают границы тонких штрихов и оставляют модели больше свободы переинтерпретации формы, усиливая геометрический дрейф; scribble сознательно эскизен и слишком свободен; lineart настроен на чистую иллюстративную графику и навязывает чуждый домену прайор. Карты глубины и нормалей неинформативны для плоского рендера, у которого нет рельефа: чтобы задействовать их, пришлось бы синтезировать карту рельефа реза, что вынесено в направления развития. Семантическая сегментация оперирует классами регионов, а не линиями, и для штриховой структуры бесполезна. Особый случай~-- MLSD: руны действительно преимущественно прямолинейны (ствол и наклонные ветви), что делает детектор прямых отрезков внешне привлекательным, однако он отбрасывает кривые элементы и знак-разделитель и навязывает прайор архитектурных линий, тем самым систематически искажая графемный инвентарь. Таким образом, Canny оказывается единственной модальностью, сохраняющей точную топологию штриха при сохранении управляемости жёсткостью обусловливания.
 
\subsubsection{Конвейер SD3 + ControlNet Canny: параметры}
\label{sec:synth-pipeline}
 
Реализованный основной конвейер устроен следующим образом. Контрастный рендер переводится в карту рёбер детектором Canny с нижним и верхним порогами~50 и~150 и однопиксельной дилатацией, утолщающей рёбра до устойчивой к шумам ширины. Полученная карта используется как управляющее изображение для модели Stable Diffusion~3 Medium~\cite{esser2024sd3} с адаптером ControlNet \texttt{InstantX/SD3-Controlnet-Canny}~\cite{zhang2023controlnet}. Под ограничение видеопамяти порядка~15,6~ГБ модель загружается в формате \texttt{float16} с выгрузкой блоков на CPU (model offload) и с отключением третьего текстового энкодера (T5), не критичного для коротких, фиксированных промптов.
 
Текстовое обусловливание задаётся неизменной парой промптов: положительный описывает глубоко врезанные руны на тёмно-сером гранитном носителе с выветренной эродированной поверхностью при ровном рассеянном освещении в стилистике археологической макросъёмки, отрицательный подавляет полированные, шрифтовые, «фэнтезийные» и иллюстративные артефакты, а также наложенный текст и водяные знаки. Из-за лимита в 77 токенов в CLIP-архитектурах используются самые основные и важные характеристики. Число шагов диффузии~-- 28, шкала guidance~-- 7,0.
 
Ключевой управляющий параметр~-- сила обусловливания ControlNet, зафиксированная на значении~0,65. Его повышение приближает форму синтезированных рун к управляющей карте, понижение даёт более убедительную каменную поверхность ценой большего геометрического дрейфа. Значение~0,65 выбрано как рабочий компромисс между фотореалистичностью врезки и сохранением геометрии знака. Каждое изображение сохраняется вместе с меткой, генерация ведётся возобновляемыми пакетами с дозаписью таблицы меток, что позволяет наращивать корпус инкрементально в пределах вычислительного бюджета.
 
\subsubsection{Точность меток и направления развития синтеза}
\label{sec:synth-fidelity}
 
Точность разметки синтетического корпуса требует точной формулировки, разводящей два разных утверждения. Метка-строка точна по построению: обратимое отображение однозначно фиксирует, что изображение обязано показывать. Открытым остаётся вопрос о визуальной точности~-- насколько верно изображение реализует эту строку,~-- и именно здесь два режима синтеза расходятся. В режиме inpainting визуальная реализация точна, поскольку зона знака не синтезируется; в режиме ControlNet Canny она приближённа: при значении силы обусловливания~0,65 подавляющее большинство генераций воспроизводит заданную строку верно, однако часть образцов демонстрирует дрейф геометрии, способный сделать один знак визуально близким к другому.
 
Из таких ограничений вытекает несколько направлений развития синтеза, способных внести вклад в сокращение синтетически-реального разрыва. Прямое снижение шума разметки достигается фильтром согласованности: прогон вспомогательного распознавателя по сгенерированным изображениям и отбраковка образцов, чьё прочтение расходится с целевой строкой, отсекают наиболее дрейфовавшие генерации.  Сокращение текстурного разрыва достижимо специализацией диффузионного прайора (DreamBooth\,/\,LoRA) на реальных фотографиях каменных надписей, не пересекающихся с тестовым gold set. 
 
\subsection{Предобработка изображений и протокол разбиения}
\label{sec:preprocessing}
 
\subsubsection{Предобработка изображений}
\label{sec:preprocessing-images}
 
Принятая стратегия предобработки минимальна и сохраняет метку: поскольку деградация рунического носителя является структурным, а не аддитивным шумом, классические методы извлечения сигнала~-- бинаризация Otsu и Sauvola, выделение границ Canny~-- не используются как основные операции. Вместо этого эффективнее дообучить энкодер. Изображение приводится лишь к входному формату конкретной модели: процессор TrOCR масштабирует его к фиксированному квадратному разрешению в трёхканальном представлении, мультимодальная модель семейства Qwen-VL использует собственный препроцессор с динамическим разрешением. Цветовое представление сохраняется трёхканальным, несмотря на близкую к монохромной природу каменных снимков, поскольку предобученные энкодеры рассчитаны на RGB-вход.
 
Аугментация применяется только к синтетической обучающей выборке и носит дополняющий характер, так как основная вариативность уже внесена самим синтезом (текстура, выветривание, освещение). Используются ограниченные фотометрические преобразования, такие как: яркость, контраст, гамма, умеренное размытие и гауссов шум, а также небольшие перспективные искажения; геометрические преобразования ограничены так, чтобы не нарушать читаемость и валидность метки. Gold set, как отражение работы конвейера в реальных условиях, не аугментируется. 
 
Поскольку этап детекции в работе не используется, первичным входом распознавателя служит изображение надписи целиком. Разметка строкового уровня позволяет при необходимости подавать на вход вырезанные области строк, что особенно существенно для протяжённых надписей с дуговым и спиральным расположением знаков, образующих заведомо трудный случай и обсуждаемых среди ограничений работы.
 
\subsubsection{Протокол разбиения и контроль утечки}
\label{sec:preprocessing-split}
 
Синтетический корпус делится на обучающую и валидационную подвыборки, причём разбиение производится по уникальным словам и фразам: лексическая непересекаемость синтетических train и val гарантирует, что валидация измеряет обобщение на не встречавшиеся строки, а не запоминание; синтетическая валидация используется для подбора гиперпараметров и ранней остановки. Реальный gold set целиком образует тестовое множество и не участвует ни в обучении, ни в подборе гиперпараметров. При необходимости честной ранней остановки из него может выделяться небольшая девелопмент-доля, не пересекающаяся с тестовой.
 
Отдельного внимания требует контроль лексической утечки между синтетическим обучением и реальным тестом. Поскольку и синтетическая лексика, и тексты gold set восходят к одним и тем же реальным базам, одно и то же слово может присутствовать в обоих, и без специальной меры оценка на gold set смешивала бы перенос на реальные носители с припоминанием уже виденной строки. Чтобы оценка отражала именно визуальный перенос, слова, входящие в gold set, исключаются из лексикона синтетического обучения.
 

\section{Методология и архитектура}
\label{sec:method}
 
Работа намеренно сопоставляет два класса распознавателей, оба в сквозном (end-to-end) виде и оба адаптируемые параметрически эффективными методами на одном и том же синтетическом корпусе: специализированную OCR-модель TrOCR и мультимодальную модель общего назначения семейства Qwen-VL. Содержательный вопрос сопоставления состоит в следующем: переносится ли широкий зрительно-языковой прайор крупной универсальной модели на чуждый рунический домен лучше, чем специализированный, но меньший по размеру OCR-распознаватель, при равных обучающих данных и едином протоколе оценки.
 
\subsection{Постановка задачи и общая схема конвейера}
\label{sec:method-setup}
 
Задача распознавания формализуется как отображение $f_\theta: \mathcal{I} \to \Sigma^{*}$, переводящее изображение надписи $I \in \mathcal{I}$ в строку транслитерации над фиксированным алфавитом $\Sigma$ из 33 знаков чистого графемно-транслитерационного ядра конвенции Rundata. Отображение реализуется целиком, без промежуточного этапа детекции и классификации отдельных графем. Качество $f_\theta$ измеряется стандартным CER (4) на двух множествах: отложенной синтетической подвыборке и реальном gold set. Первое характеризует обучаемость на синтетическом распределении, второе — переносимость на реальные носители; их разность оценивает синтетически-реальный разрыв.
 
Принципиально, что для обеих архитектур целевой выход~-- это латинская строка транслитерации, а не последовательность кодовых позиций рунического блока Unicode. Это решение имеет прямое архитектурное следствие, общее для TrOCR и Qwen-VL: алфавит $\Sigma$ целиком покрывается уже существующими субсловными словарями предобученных моделей, что устраняет необходимость расширять токенизатор рунами и позволяет переиспользовать предобученные текстовые представления без модификации~-- критическое преимущество в условиях малого объёма данных. Сводное сопоставление двух распознавателей приведено в таблице~\ref{tab:baselines}.

Переход к сквозному предсказанию строки снимает необходимость в явной system-aware голове классификатора для покомпонентной архитектуры: разрешение системной принадлежности и графемно-фонемная неоднозначность перекладываются на авторегрессивный языковой декодер, моделирующий целевую строку контекстно. Тем самым требование системной обусловленности не отменяется, а реализуется неявно — ценой потери интерпретируемости пографемной классификации, что обсуждается среди ограничений
 
\begin{table}[htbp]
\centering
\caption{Сопоставление базовых линий распознавания}
\label{tab:baselines}
\small
\begin{tabular}{p{3.0cm} p{4.6cm} p{4.6cm}}
\hline
\textbf{Свойство} & \textbf{TrOCR} & \textbf{Qwen-VL} \\
\hline
Класс модели & специализированный OCR-распознаватель & мультимодальная модель общего назначения \\
Энкодер & визуальный трансформер (ViT\,/\,BEiT\,/\,DeiT) & визуальный трансформер с нативным динамическим разрешением \\
Декодер & текстовый трансформер (RoBERTa\,/\,MiniLM) & языковая модель Qwen (авторегрессивная) \\
Обработка входа & масштабирование к фиксированному квадрату & переменное разрешение, инструкция + изображение \\
Адаптация & LoRA (база заморожена) & QLoRA: 4-битная база + LoRA \\
Целевой выход & строка транслитерации (алфавит $\Sigma$) & строка транслитерации (алфавит $\Sigma$) \\
Испытанные размеры & базовый уровень (BASE) & $\approx$\,3 и $\approx$\,7 млрд параметров; далее Qwen3-VL \\
Роль в работе & специализированная базовая линия & универсальная базовая линия \\
\hline
\end{tabular}
\end{table}
 
 
\subsection{Baseline 1: TrOCR с адаптацией LoRA}
\label{sec:method-trocr}
 
Первый baseline~-- модель TrOCR~\cite{li2023}, объединяющая визуальный трансформерный энкодер (ViT\,/\,BEiT\,/\,DeiT) с авторегрессивным текстовым декодером (RoBERTa\,/\,MiniLM) и предобученная на масштабном двухступенчатом корпусе печатных и рукописных строк. Выбор TrOCR как специализированной линии мотивирован тем, что в литературе по историческому HTR она образует фактический стандарт и, что существеннее для малоресурсной задачи, допускает дообучение относительно небольшим числом размеченных строк~\cite{strobel2022,vogler2022}. При этом сохраняется зафиксированная в обзоре оговорка: устойчивое превосходство над свёрточно-рекуррентными конвейерами демонстрирует прежде всего вариант LARGE~\cite{strobel2023}, тогда как меньшие конфигурации могут уступать. В работе под ограничение видеопамяти используется базовая конфигурация (BASE), и расхождение с LARGE при наличии бюджета фиксируется как фактор сравнения.
 
Адаптация выполняется методом LoRA~\cite{hu2022}: в линейные проекции слоёв внимания и сопутствующие полносвязные подслои внедряются обучаемые низкоранговые матрицы, тогда как исходные веса остаются замороженными. Это сокращает число обучаемых параметров на порядки, что одновременно снижает требования к видеопамяти и ограничивает переобучение на синтетическом распределении. Существенно, что адаптеры охватывают не только энкодер, но и декодер: как показано Parres и Paredes~\cite{parres2023}, при адаптации к новому языковому материалу замораживание декодера ведёт к потере нескольких процентных пунктов CER, поэтому адаптация только зрительной части в данной задаче недостаточна. Целевой алфавит $\Sigma$ является подмножеством предобученного субсловного словаря, благодаря чему дообучение не требует расширения токенизатора. Обучение ведётся на синтетическом корпусе, отбор контрольной точки и ранняя остановка~-- по синтетической валидации с лексической непересекаемостью.
 
\subsection{Baseline 2: мультимодальная модель Qwen-VL с адаптацией QLoRA}
\label{sec:method-qwen}
 
Второй baseline~-- мультимодальная модель семейства Qwen-VL~\cite{wang2024qwen2vl}, в которой визуальный энкодер преобразует изображение в последовательность зрительных токенов, подаваемых вместе с текстовой инструкцией в авторегрессивную языковую модель, генерирующую ответ. Архитектурно значимы две особенности. Во-первых, нативная поддержка переменного разрешения позволяет подавать изображение надписи без агрессивного масштабирования к фиксированному квадрату, что потенциально сохраняет мелкие детали штриха, критичные для различения аллографов. Во-вторых, задача задаётся инструкцией: модель получает изображение и предписание выдать транслитерацию в конвенции Rundata, а целевым выходом, как и для TrOCR, служит строка над алфавитом $\Sigma$.
 
Использование универсальной мультимодальной модели как распознавателя требует обоснования, поскольку систематические оценки показывают, что в режиме zero-shot такие модели массово проваливаются именно на редких скриптах, мелких символах и нестандартной композиции~\cite{liu2024ocrbench,fu2025ocrbench2}~-- всех трёх характеристиках рунической эпиграфики. Отсюда следует, что универсальная модель пригодна здесь не как готовый распознаватель, а как объект дообучения, и сопоставление с TrOCR ведётся именно в режиме адаптации. Адаптация выполняется методом QLoRA: базовая модель квантуется в 4-битное представление, а обучаемые низкоранговые адаптеры LoRA внедряются в проекции внимания языковой части. Квантование базовых весов снижает объём занимаемой видеопамяти настолько, что модель с семью миллиардами параметров умещается в бюджет порядка~15,6~ГБ, оставаясь обучаемой через адаптеры.
 
Ограниченность видеопамяти задаёт единый принцип отбора моделей. Для каждого семейства выбирается наибольший вариант, чья пиковая память обучения не превосходит бюджета при фиксированных квантовании, ранге адаптеров, размере пакета и разрешении входа. Сообщаются все варианты, удовлетворяющие этому ограничению. В соответствии с этим критерием в работе испытаны конфигурации Qwen2.5-VL 3B, 7B, а также Qwen3-VL 3B, 7B.
 
\subsection{Демонстрационный модуль лингвистической интерпретации}
\label{sec:method-interp}

 
\subsection{Протокол обучения и инфраструктура}
\label{sec:method-infra}
 
Все эксперименты выполняются на единственном графическом ускорителе с объёмом видеопамяти порядка~15,6~ГБ в средах Colab\,/\,Kaggle, что и определяет выбранные решения. Вычисления ведутся в смешанной точности (bf16\,/\,fp16); базовые веса Qwen-VL хранятся в 4-битном представлении (NF4). Конфигурация LoRA~-- ранг, масштабирующий коэффициент, dropout и множество целевых модулей~-- приводится в репрезентативных пределах, а её итоговые значения, как и значения скорости обучения, числа эпох и эффективного размера пакета фиксируются в экспериментальной части. Оптимизатор~-- AdamW; ранняя остановка и отбор контрольной точки~-- по синтетической валидации. Прогон каждой конфигурации сопровождается логированием гиперпараметров, кривых обучения и метрик в систему трекинга экспериментов (Weights~\&~Biases\) и фиксацией начальных состояний генераторов случайных чисел ради воспроизводимости. Единообразие протокола между TrOCR и Qwen-VL~-- один и тот же обучающий корпус, одна и та же синтетическая валидация, один и тот же тестовый gold set и единая метрика CER~-- обеспечивает корректность сравнения двух классов моделей.
\newpage 
\printbibliography[heading=bibintoc] 

% \begin{thebibliography}{0}
% 	\bibitem{chirkova18}\hypertarget{chirkova18}{}
% 	\href{https://arxiv.org/abs/1810.10927}
% 	{Nadezhda Chirkova, Ekaterina Lobacheva, Dmitry Vetrov. Bayesian Compression for Natural Language Processing. In EMNLP 2018.}
% \end{thebibliography}
	
\newpage
\appendix

\section{Пример секции аппендикса}

\end{document}
